% Generated draft: requires mathematical and source review.
\documentclass[11pt]{article}
\usepackage[margin=25mm]{geometry}
\usepackage{fontspec}
\setmainfont{Latin Modern Roman}
\usepackage{amsmath,amssymb,mathtools,unicode-math}
\setmathfont{Latin Modern Math}
\usepackage{enumitem,needspace,xurl,hyperref}
\setlength{\parindent}{0pt}
\setlength{\parskip}{5pt}
\setlength{\emergencystretch}{3em}
\newcommand{\N}{\mathbb{N}}
\newcommand{\Z}{\mathbb{Z}}
\newcommand{\Q}{\mathbb{Q}}
\newcommand{\R}{\mathbb{R}}
\newcommand{\C}{\mathbb{C}}
\newcommand{\F}{\mathbb{F}}
\newcommand{\A}{\mathbb{A}}
\newcommand{\PP}{\mathbb{P}}
\newcommand{\E}{\mathbb{E}}
\newcommand{\eps}{\varepsilon}
\newcommand{\dd}{\,\mathrm{d}}
\begin{document}
\section*{UnsolvedMath ID: 30005097}
\textbf{Source problem code:} catalogue entry
\subsection{Definitions and mathematical statement}
Let \(K\) be a number field, that is, a finite extension of \(\Q\), and put \(n=[K:\Q]\). The supplied question concerns a set denoted by \(\operatorname{PI}(K)\), with Hermite equivalence and \(\operatorname{GL}_2(\Z)\)-equivalence. The record does not define this set or either equivalence relation. Consequently, these are unresolved source-dependent notions, and the formulation below is conditional on their intended definitions.

For \(R\in\{\Z,\Q\}\), let \(\operatorname{GL}_2(R)\) denote the group of invertible \(2\times2\) matrices over \(R\). Specifying this group does not specify the action or equivalence relation intended in the source. Write \(\sim_H\) for the intended Hermite equivalence and \(\sim_{\Z}\) for the intended \(\operatorname{GL}_2(\Z)\)-equivalence on \(\operatorname{PI}(K)\). The wording of the question presupposes that each Hermite-equivalence class is a union of \(\sim_{\Z}\)-classes; the missing definitions must establish this compatibility.

Subject to these qualifications, define
\[
\mathcal S(K)=\left\{H\in\operatorname{PI}(K)/{\sim_H}:\ \exists f,g\in H\text{ with }f\not\sim_{\Z}g\right\}.
\]
The intended question is whether
\[
\text{for every number field }K,\qquad |\mathcal S(K)|<\infty.
\]
Thus, for each fixed \(K\), are there only finitely many Hermite-equivalence classes containing more than one \(\operatorname{GL}_2(\Z)\)-equivalence class? The proposed finiteness is for each field separately; no bound uniform in \(K\) is supplied.

The original record also supplies the following contextual assertions, which are retained as source assertions rather than independently verified theorems:
\begin{itemize}
\item Theorem 2, attributed there to reference [4]: for every integer \(n\geq4\), \(\operatorname{PI}(n)\) contains infinitely many Hermite-equivalence classes that are unions of at least two \(\operatorname{GL}_2(\Z)\)-equivalence classes. Neither \(\operatorname{PI}(n)\) nor reference [4] is identified in the supplied record. Its relationship to the fixed-field set \(\operatorname{PI}(K)\) is therefore unresolved.
\item Theorem 3: if \(K\) has degree \(n\geq5\) and its Galois closure has Galois group \(S_n\), then the Hermite-equivalence classes in \(\operatorname{PI}(K)\) that split into more than one \(\operatorname{GL}_2(\Z)\)-equivalence class lie in finitely many \(\operatorname{GL}_2(\Q)\)-equivalence classes. Here \(S_n\) is the symmetric group on \(n\) letters. The intended rational equivalence relation is also undefined in the record. This supplied assertion concerns finiteness of rational equivalence classes, which does not by itself state the requested finiteness of Hermite-equivalence classes.
\end{itemize}
\subsection{Short English statement}
For each fixed number field \(K\), ask whether only finitely many Hermite-equivalence classes in \(\operatorname{PI}(K)\) split into multiple \(\operatorname{GL}_2(\Z)\)-equivalence classes. The supplied record omits the definitions of \(\operatorname{PI}(K)\) and the equivalence relations, so this is a conditional statement-only draft requiring those definitions and human mathematical review.
\subsection{Sources}
\begin{itemize}
\item Local UnsolvedMath registry, supplied record with permanent ID \(30005097\), title ``Hermite and GL2 Equivalence Classes of Number Fields,'' exact source problem code \texttt{catalogue entry}, in the Oberwolfach Reports Open Problems collection. The proposed year is supplied as 2022; no proposer is named.
\item Supplied report citation: \emph{Diophantische Approximationen} (2023), Oberwolfach Reports, pp. 1091--1164, DOI \texttt{10.4171/owr/2022/21}: \url{https://doi.org/10.4171/owr/2022/21}. The literature-source label for this link is ``Evertse et al., Hermite equivalence.'' The differing year indications are retained as supplied.
\item Supplied literature reference: ``Bhargava et al., Hermite equivalence of polynomials'': \url{https://arxiv.org/abs/2109.02932}.
\item The original statement attributes its Theorem 2 to reference [4], whose bibliographic identity is absent from the supplied record. No identification with either listed reference is assumed.
\item Literature assessment checked at \texttt{2026-08-22}: the supplied status is ``partially solved.'' The supplied assessment reports finite-class results for suitable symmetric Galois closures and infinitely many split Hermite classes, while describing the exact finiteness question for arbitrary number fields as open. This is supplied dated curation derived from the OWR material, not independently verified here; the record itself requests independent verification of time-sensitive status and cited claims.
\item The registry labels the statement \texttt{corrected\_verified} and says it was checked against the official Oberwolfach report. These are supplied verification claims, not a claim that the present draft independently consulted or verified that report.
\end{itemize}
\end{document}
